% Figure: a floating iceberg with about nine-tenths of its volume submerged,
% illustrating the density-ratio buoyancy balance. No em-dashes.
\begin{figure}[htbp]
\centering
\begin{tikzpicture}[scale=1.0]
% water and air regions
\fill[engblue!12] (-0.5,0) rectangle (9.5,3.0);
\draw[engblue, thick] (-0.5,3.0) -- (9.5,3.0);
\node[engblue, font=\scriptsize, anchor=south] at (-0.3,3.0) {waterline};
% iceberg: an irregular polygon crossing the waterline
% above-water tip (small), below-water bulk (large)
\fill[engamber!30, draw=engblack, thick]
  (3.6,3.9) -- (4.9,3.55) -- (5.6,3.2) -- (5.9,2.4) -- (5.5,1.2)
  -- (4.6,0.5) -- (3.4,0.8) -- (3.0,2.0) -- (3.1,3.0) -- (3.6,3.9) -- cycle;
% waterline crossing shading of submerged part
\node[font=\scriptsize] at (4.3,3.5) {tip};
\node[font=\scriptsize] at (4.5,1.6) {submerged bulk};
% fraction labels
\draw[enggray, dashed] (6.1,3.9) -- (7.4,3.9);
\draw[enggray, dashed] (6.1,3.0) -- (7.4,3.0);
\draw[enggray, dashed] (6.1,0.5) -- (7.4,0.5);
\draw[engred, <->] (7.2,3.9) -- (7.2,3.0);
\node[engred, font=\scriptsize, anchor=west] at (7.2,3.45) {$\sim 0.1$ above};
\draw[engblue, <->] (7.2,3.0) -- (7.2,0.5);
\node[engblue, font=\scriptsize, anchor=west] at (7.2,1.75) {$\sim 0.9$ below};
% force arrows
\draw[engblack, ultra thick, ->] (4.4,2.2) -- (4.4,1.0);
\node[engblack, font=\scriptsize, anchor=west] at (4.45,1.5) {$W$};
\draw[englime!70!black, ultra thick, ->] (4.9,1.8) -- (4.9,3.0);
\node[englime!60!black, font=\scriptsize, anchor=west] at (4.95,2.5) {$F_B$};
\end{tikzpicture}
\caption[Floating fraction of an iceberg]{An iceberg floats with the
submerged fraction equal to the density ratio $\rho_{\mathrm{ice}}/
\rho_{\mathrm{sea}}$. Glacial ice at about $917\,\mathrm{kg/m^3}$ in
seawater at about $1025\,\mathrm{kg/m^3}$ leaves roughly nine-tenths of
the volume below the waterline and one-tenth showing, the origin of the
tip-of-the-iceberg image. The weight $W$ acting down through the centre
of mass balances the buoyancy $F_B$ acting up through the centroid of the
displaced volume. When floating ice melts it releases exactly the volume
of water it was already displacing, so the level of the surrounding water
does not change, the reasoning worked in the text.}
\label{fig:vol03ch10:iceberg}
\end{figure}
